\documentclass{situationist-preprint}

\title{Republican Freedom as Non-Domination}
\msid{POL-09-b}
\msfamily{POL-09}
\msdiscipline{Political Philosophy}
\mstype{Consequence}
\mskeywords{republicanism; non-domination; Pettit; negative liberty; positive liberty; robustness}

\begin{document}
\maketitle

\begin{abstract}
Republicans hold that freedom is non-domination: not being subject to the capacity of another to interfere arbitrarily in one's choices. They present it as a third concept, neither negative nor positive. POL-09 concludes that freedom has exactly two irreducible dimensions, the negative and the positive. We derive what follows for non-domination. It is not a third dimension but a security condition on the negative one: it requires that non-prevention be robust across variations in the will of others. What republicans add beyond this, the servility and self-censorship that domination breeds, belongs to the positive dimension. The republican's cases are thereby explained, and the two-dimensional account is preserved.
\end{abstract}

\section{The republican challenge}

\citet{Pettit1997} and \citet{Skinner1998} argue that the tradition of the free state understood freedom as the condition of a person who is not subject to the arbitrary will of another. The slave of a benevolent master, who never interferes, is not interfered with, and yet he is unfree; the citizen of a well-ordered republic is interfered with by law, and yet, if the law is non-arbitrary, she is free. If both verdicts are correct, freedom seems to be neither absence of interference nor self-mastery, but a third thing.

POL-09-a argues that freedom has a negative dimension, the extent to which no other agent prevents a person from pursuing an option, and a positive dimension, the extent to which she has the capacities to direct her own pursuit of it (Definition~2.1 of POL-09-a), and that neither reduces to the other (Theorem~3.2 of POL-09-a). The republican challenge is that this leaves out a third dimension.

\section{Non-domination as a security condition}

The negative dimension measures the \emph{extent} to which no one prevents a person from pursuing an option. Extent has a modal measure as well as an actual one: a person whose non-prevention holds only because another happens not to want to prevent her enjoys it in fewer of the relevant possibilities than a person whose non-prevention holds whatever others want.

\begin{definition}\label{d:robust}
A person's non-prevention with respect to $z$ is \emph{robust} to the extent that it would persist under variations in the will of other agents.
\end{definition}

\begin{proposition}\label{p:security}
The benevolent master's slave and the free citizen differ in the negative dimension, read with the robustness of Definition~\ref{d:robust}.
\end{proposition}

\begin{proof}
The slave's non-prevention depends on the master's will: were the master to change his mind, the slave would be prevented, and nothing stands in the way. The citizen's non-prevention in the range of options the law protects persists whatever any particular agent wills, since the law constrains them. So the citizen's non-prevention is robust and the slave's is not.
\end{proof}

Proposition~\ref{p:security} shows that the central republican case can be described without a third dimension: the slave lacks secure non-prevention. This is the structure of the account of freedom as independence of \citet{List2016}, who argue that freedom so understood is robust non-interference. The citizen's case is equally accommodated. Interference by law is prevention, and it limits the negative dimension in the options it forecloses; but the law that forecloses some options secures non-prevention in many more, and the republican verdict that she is free concerns that larger range.

\section{The residue belongs to the positive dimension}

Republicans claim more. The dominated person, even when never interfered with, is diminished: she anticipates the master's wishes, censors herself, and adopts the deference of one who lives at another's mercy \citep{Pettit1997}. These are not facts about prevention, actual or counterfactual. They are deficits in the capacity to form, assess and act on one's own ends, and so belong to the positive dimension of Definition~2.1 of POL-09-a. The same holds of a related question, raised by the analysis of power in POL-10: whether a preference shaped by the power of others is the person's own is a question about the positive dimension of her freedom.

\begin{corollary}\label{c:supervene}
Any two persons alike in robust non-prevention and in the capacities for self-direction are alike in non-domination.
\end{corollary}

\begin{proof}
By Proposition~\ref{p:security} and the argument of this section, differences in domination appear either as differences in the robustness of non-prevention or as differences in the capacities that domination erodes. If neither differs, nothing remains in which a difference in domination could consist.
\end{proof}

\begin{objection}
Robust non-prevention is not non-domination: a powerful neighbor who could invade but is checked by mutual deterrence leaves non-prevention fairly robust, yet he dominates.
\end{objection}

\begin{reply}
If he could invade at will, the non-prevention is not robust with respect to his will; if he is checked, he does not dominate in Pettit's sense, since his capacity to interfere is not arbitrary. The case does not separate the two. What it shows is that robustness is a matter of degree, which is what the negative dimension, as an extent, allows.
\end{reply}

\section{Conclusion}

Republican freedom is real and important, but it is not a third dimension. Non-domination is a security condition on the negative dimension, and the further harms of domination are harms to the positive dimension. The two-dimensional account of POL-09 thus explains the republican's verdicts without adding to its dimensions. The further question republicans press, whether a capacity to interfere is checked and answerable, is the question that POL-10-a assigns to the theory of domination rather than to the concept of power.

\begin{thebibliography}{9}
\bibitem[List and Valentini(2016)]{List2016} List, C. and Valentini, L. (2016). Freedom as independence. \emph{Ethics}, 126(4), 1043--1074.
\bibitem[Pettit(1997)]{Pettit1997} Pettit, P. (1997). \emph{Republicanism: A Theory of Freedom and Government}. Oxford University Press.
\bibitem[Skinner(1998)]{Skinner1998} Skinner, Q. (1998). \emph{Liberty before Liberalism}. Cambridge University Press.
\end{thebibliography}

\end{document}
